
We can now rewrite the primal problem~\eqref{eq:primal-problem} as:
\begin{align*}
	\min_{\symbf{w}, b, \symbf{\xi}} p(\symbf{w}, b, \symbf{\xi}) & = \min_{\symbf{w}, b, \symbf{\xi}} \Lcal(\symbf{w}, b, \symbf{\xi}, \symbf{\alpha}_1, \symbf{\alpha}_2)                                                                                                        \\
	                                                              & = \min_{\symbf{w}, b, \symbf{\xi}} \max_{\symbf{\alpha}_1, \symbf{\alpha}_2} \Lcal(\symbf{w}, b, \symbf{\xi}, \symbf{\alpha}_1, \symbf{\alpha}_2)                                                              \\
	                                                              & = -\frac{1}{2}\inner*{\symbf{\alpha}_1, \mathrm{Y} \mathrm{G} \mathrm{Y} \symbf{\alpha}_1} +  \inner*{\symbf{1}_M, \symbf{\alpha}_1} + C\symbf{1}_M^\top \symbf{\xi} -  \inner*{\symbf{0}_M, \symbf{\alpha}_2}
\end{align*}

\begin{equation}
	\min_{\symbf{w}, b, \symbf{\xi}} p(\symbf{w}, b, \symbf{\xi}) = \min_{\symbf{w}, b, \symbf{\xi}} \max_{\symbf{\alpha}_1, \symbf{\alpha}_2} \Lcal(\symbf{w}, b, \symbf{\xi}, \symbf{\alpha}_1, \symbf{\alpha}_2)
\end{equation}

where we now define:

\[
	q(\symbf{\alpha}_1,\symbf{\alpha}_2) = \min_{\symbf{w}, b, \symbf{\xi}} \Lcal(\symbf{w}, b, \symbf{\xi}, \symbf{\alpha}_1, \symbf{\alpha}_2) \tag{D}
\]

Thus, we can define the \emph{Lagrangian dual} of the primal problem~\eqref{eq:primal-problem} as:
\begin{align*}
	\max_{\symbf{\alpha}_1, \symbf{\alpha}_2} q(\symbf{\alpha}_1,\symbf{\alpha}_2) & = \max_{\symbf{\alpha}_1, \symbf{\alpha}_2} \min_{\symbf{w}, b, \symbf{\xi}} \Lcal(\symbf{w}, b, \symbf{\xi}, \symbf{\alpha}_1, \symbf{\alpha}_2)                                            \\
	                                                                               & = \max_{\symbf{\alpha}_1, \symbf{\alpha}_2} -\left(\frac{1}{2}\inner*{\symbf{\alpha}_1, \mathrm{Y} \mathrm{G} \mathrm{Y} \symbf{\alpha}_1} -  \inner*{\symbf{1}_M, \symbf{\alpha}_1} \right) \\
\end{align*}

Now, if we want to minimize $q$ instead of maximizing it, we can simply multiply the entire equation by -1, which gives us:
\begin{align*}
	\min_{\symbf{\alpha}_1, \symbf{\alpha}_2} q(\symbf{\alpha}_1,\symbf{\alpha}_2) & = -\max_{\symbf{\alpha}_1, \symbf{\alpha}_2} -q(\symbf{\alpha}_1,\symbf{\alpha}_2)                                                                                                                                             \\
	                                                                               & = -\max_{\symbf{\alpha}_1, \symbf{\alpha}_2} \min_{\symbf{w}, b, \symbf{\xi}} -\left(\frac{1}{2}\inner*{\symbf{\alpha}_1, \mathrm{Y} \mathrm{G} \mathrm{Y} \symbf{\alpha}_1} -  \inner*{\symbf{1}_M, \symbf{\alpha}_1} \right) \\
	                                                                               & = \min_{\symbf{\alpha}_1, \symbf{\alpha}_2} \left(\frac{1}{2}\inner*{\symbf{\alpha}_1, \mathrm{Y} \mathrm{G} \mathrm{Y} \symbf{\alpha}_1} -  \inner*{\symbf{1}_M, \symbf{\alpha}_1} \right)                                    \\
	                                                                               & = \min_{\symbf{\alpha}_1, \symbf{\alpha}_2} -\left(\frac{1}{2}\inner*{\symbf{\alpha}_1, \mathrm{Y} \mathrm{G} \mathrm{Y} \symbf{\alpha}_1} -  \inner*{\symbf{1}_M, \symbf{\alpha}_1} \right)                                   \\
\end{align*}

This gives us the dual problem, where \(\symbf{\alpha} := \symbf{\alpha}_1\) and \(\symbf{\alpha}_2 = \symbf{0}\):

\begin{align}\label{eq:dual-problem}
	\min_{\symbf{\alpha}} & \frac{1}{2} \inner*{\symbf{\alpha}, \mathrm{Y} \mathrm{G} \mathrm{Y} \symbf{\alpha}} -  \inner*{\symbf{1}_M, \symbf{\alpha}}
	\quad
	\text{subject to}
	\quad
	\begin{cases}
		0 \leq \symbf{\alpha} \leq C\symbf{1}_M \\
		\symbf{\alpha}^\top Y = 0
	\end{cases}\tag{DP}
\end{align}\qedhere



We derive the dual problem from the primal problem~\eqref{eq:primal-problem} using Lagrange multipliers and the KKT conditions.

Here we show how to derive the dual problem ~\eqref{eq:dual-problem} from the primal problem \eqref{eq:primal-problem} using Lagrange multipliers.

We start with the primal problem~\eqref{eq:primal-problem}:
\begin{align}
	\min_{\symbf{w}, b, \symbf{\xi}} & \frac{1}{2} \norm*{\symbf{w}}_2^2 + C\symbf{1}_M^\top \symbf{\xi}
	\quad \text{subject to} \quad
	\begin{cases}
		\symbf{1}_M - Y(X\symbf{w} + b\symbf{1}_M) - \symbf{\xi} \leq \symbf{0}_M \\
		-\symbf{\xi} \leq \symbf{0}_M
	\end{cases}
\end{align}

We want to solve the minimization problem~\eqref{eq:primal-problem} with the constraints \(c_i(\symbf{w}, b, \symbf{\xi}) \leq 0\).
An equivalent, but not very efficient, way to do this is to minimize the function:

\begin{align*}\label{eq:penalty-function}
	p(\symbf{w}, b, \symbf{\xi}) & =
	\begin{cases}
		f(\symbf{w}, b, \symbf{\xi}) & \text{if } c_i(\symbf{w}, b, \symbf{\xi}) \leq 0 \\
		+\infty                      & \text{otherwise}
	\end{cases}                             \\
	                             & = f(\symbf{w}, b, \symbf{\xi}) + \mathcal{I}(c_i(\symbf{w}, b, \symbf{\xi})) \\
\end{align*}
where \(\mathcal{I}\) is the indicator function, which is defined as:
\[
	\mathcal{I}(c_i(\symbf{w}, b, \symbf{\xi})) =
	\begin{cases}
		0       & \text{if } c_i(\symbf{w}, b, \symbf{\xi}) \leq 0 \\
		+\infty & \text{otherwise}
	\end{cases}
\]

Now this is not a very efficient way to solve the problem, since \(\mathcal{I}\) is both non-differentiable and discontinuous.

\paragraph{Lagrangian}
Instead, we can use the Lagrangian function to solve the problem. The Lagrangian function is a way to incorporate the constraints into the objective function by introducing Lagrange multipliers.
Let \(\symbf{\alpha}_1, \symbf{\alpha}_2 \geq 0\) be the Lagrange multipliers for the constraints \(c_i(\symbf{w}, b, \symbf{\xi}) \leq 0\).

Then we can write the Lagrangian function as:

\begin{align*}
	\Lcal(\symbf{w}, b, \symbf{\xi}, \symbf{\alpha}_1, \symbf{\alpha}_2) & = f(\symbf{w}, b, \symbf{\xi}) + \symbf{\alpha}_1^\top c_1(\symbf{w}, b, \symbf{\xi}) + \symbf{\alpha}_2^\top c_2(\symbf{\xi})
\end{align*}

If the constraints \(c_i(\symbf{w}, b, \symbf{\xi}) \leq 0\) now hold, then we can just set \(\symbf{\alpha}_1, \symbf{\alpha}_2 = 0\) and function becomes:
\[
	\Lcal(\symbf{w}, b, \symbf{\xi}, \symbf{0}, \symbf{0}) = f(\symbf{w}, b, \symbf{\xi})
\]

However, if the constraints do not hold, i.e. \(c_i(\symbf{w}, b, \symbf{\xi}) > 0\), then we can set \(\symbf{\alpha}_1, \symbf{\alpha}_2\) to be large enough so that the Lagrangian function becomes:
\[
	\Lcal(\symbf{w}, b, \symbf{\xi}, \symbf{\alpha}_1, \symbf{\alpha}_2) = f(\symbf{w}, b, \symbf{\xi}) + \infty \to \infty
\]

It so happens that if we take the \emph{maximum} of the Lagrangian with respect to \(\symbf{\alpha}_1, \symbf{\alpha}_2\), we recover the original function \(p(\symbf{w}, b, \symbf{\xi})\).

This is because, if the constraints hold, then the Lagrangian is equal to the original function \(f(\symbf{w}, b, \symbf{\xi}\), and if the constraints do not hold, then the Lagrangian is equal to \(\infty\).

\begin{align*}
	\max_{\symbf{\alpha}_1, \symbf{\alpha}_2} \Lcal(\symbf{w}, b, \symbf{\xi}, \symbf{\alpha}_1, \symbf{\alpha}_2) & = p(\symbf{w}, b, \symbf{\xi})
\end{align*}

Recall, that we still want to minimize the function \(p(\symbf{w}, b, \symbf{\xi})\) with respect to \(\symbf{w}, b, \symbf{\xi}\), so we can write the problem as:

\[
	\min_{\symbf{w}, b, \symbf{\xi}} p(\symbf{w}, b, \symbf{\xi}) = \min_{\symbf{w}, b, \symbf{\xi}} \max_{\symbf{\alpha}_1, \symbf{\alpha}_2} \Lcal(\symbf{w}, b, \symbf{\xi}, \symbf{\alpha}_1, \symbf{\alpha}_2)
\]
Since the Lagrangian \(\Lcal(\symbf{w}, b, \symbf{\xi}, \symbf{\alpha}_1, \symbf{\alpha}_2)\) is convex with respect to the primal variables \(\symbf{w}, b, \symbf{\xi}\),
and concave with respect to the dual variables \(\symbf{\alpha}_1, \symbf{\alpha}_2\),
we can apply the minimax theorem \cite{vonNeumann1928}.

\[
	\min_{\symbf{w}, b, \symbf{\xi}} \max_{\symbf{\alpha}_1, \symbf{\alpha}_2} \Lcal(\symbf{w}, b, \symbf{\xi}, \symbf{\alpha}_1,\symbf{\alpha}_2) = \max_{\symbf{\alpha}_1,\symbf{\alpha}_2} g(\symbf{\alpha}_1,\symbf{\alpha}_2)
\]

where we define:

\[
	g(\symbf{\alpha}_1,\symbf{\alpha}_2) = \min_{\symbf{w}, b, \symbf{\xi}} \Lcal(\symbf{w}, b, \symbf{\xi}, \symbf{\alpha}_1,\symbf{\alpha}_2)
\]

We can now rewrite the primal problem~\eqref{eq:primal-problem} as:
\begin{align*}
	\Lcal(\symbf{w}, b, \symbf{\xi}, \symbf{\alpha}_1, \symbf{\alpha}_2) & = \frac{1}{2} \norm*{\symbf{w}}_2^2 + C\symbf{1}_M^\top \symbf{\xi} + \symbf{\alpha}_1^\top c_1(\symbf{w}, b, \symbf{\xi}) + \symbf{\alpha}_2^\top c_2(\symbf{w}, b, \symbf{\xi})                                                                            \\
	                                                                     & = \frac{1}{2} \norm*{\symbf{w}}_2^2 + C\symbf{1}_M^\top \symbf{\xi} + \symbf{\alpha}_1^\top \left(\symbf{1}_M - Y(X\symbf{w} + b\symbf{1}_M) - \symbf{\xi}\right) - \symbf{\alpha}_2^\top \symbf{\xi}                                                        \\
	                                                                     & = \frac{1}{2} \norm*{\symbf{w}}_2^2 + C\symbf{1}_M^\top \symbf{\xi} + \symbf{\alpha}_1^\top \symbf{1}_M - \symbf{\alpha}_1^\top Y X \symbf{w} - \symbf{\alpha}_1^\top Y b\symbf{1}_M - \symbf{\alpha}_1^\top \symbf{\xi} - \symbf{\alpha}_2^\top \symbf{\xi} \\
\end{align*}

\paragraph{KKT conditions}
We can now take the gradient of the Lagrangian with respect to the primal variables \(\symbf{w}, b, \symbf{\xi}\) and set it to zero to find the optimal solution.
The gradients are given by:

Now we get that:


Substituting the KKT conditions into the Lagrangian gives us:

\begin{align*}
	\Lcal(\symbf{w}, b, \symbf{\xi}, \symbf{\alpha}_1, \symbf{\alpha}_2) & = \frac{1}{2} \norm*{\symbf{w}}_2^2 + C\symbf{1}_M^\top \symbf{\xi} + \symbf{\alpha}_1^\top \symbf{1}_M - \symbf{\alpha}_1^\top Y X \symbf{w} - \symbf{\alpha}_1^\top Y b\symbf{1}_M - \symbf{\alpha}_1^\top \symbf{\xi} - \symbf{\alpha}_2^\top \symbf{\xi}                                                                                                                                                        \\
	                                                                     & = \frac{1}{2}\inner*{\symbf{\alpha}_1, \mathrm{Y} \mathrm{G} \mathrm{Y} \symbf{\alpha}_1} + C\symbf{1}_M^\top \symbf{\xi} + \symbf{\alpha}_1^\top \symbf{1}_M - \symbf{\alpha}_1^\top Y \overbrace{X X^\top}^{G} Y \symbf{\alpha}_1 - \overbrace{\symbf{\alpha}_1^\top Y}^{\symbf{0}_M^\top} b\symbf{1}_M - \symbf{\alpha}_1^\top \symbf{\xi} - \left(C \symbf{1}_M^\top - \symbf{\alpha}_1^\top\right) \symbf{\xi} \\
	                                                                     & = \frac{1}{2}\inner*{\symbf{\alpha}_1, \mathrm{Y} \mathrm{G} \mathrm{Y} \symbf{\alpha}_1} + C\symbf{1}_M^\top \symbf{\xi} + \symbf{\alpha}_1^\top \symbf{1}_M - \symbf{\alpha}_1^\top Y G Y \symbf{\alpha}_1 -  \symbf{\alpha}_1^\top \symbf{\xi} - C \symbf{1}_M^\top \symbf{\xi} + \symbf{\alpha}_1^\top \symbf{\xi}                                                                                              \\
	                                                                     & = \frac{1}{2}\inner*{\symbf{\alpha}_1, \mathrm{Y} \mathrm{G} \mathrm{Y} \symbf{\alpha}_1} -  \inner*{\symbf{\alpha}_1, \mathrm{Y} \mathrm{G} \mathrm{Y} \symbf{\alpha}_1} + \symbf{\alpha}_1^\top \symbf{1}_M                                                                                                                                                                                                       \\
	                                                                     & = -\frac{1}{2}\inner*{\symbf{\alpha}_1, \mathrm{Y} \mathrm{G} \mathrm{Y} \symbf{\alpha}_1} +  \symbf{\alpha}_1^\top \symbf{1}_M                                                                                                                                                                                                                                                                                     \\
	\Lcal(\symbf{w}, b, \symbf{\xi}, \symbf{\alpha}_1)                   & = -\frac{1}{2}\inner*{\symbf{\alpha}_1, \mathrm{Y} \mathrm{G} \mathrm{Y} \symbf{\alpha}_1} +  \inner*{\symbf{1}_M, \symbf{\alpha}_1}
\end{align*}

If we now want to minimize \(g(\symbf{\alpha}_1,\symbf{\alpha}_2)\) instead of maximizing it, we can simply multiply the entire equation by -1, which gives us:
\begin{align*}
	\min_{\symbf{\alpha}_1, \symbf{\alpha}_2} g(\symbf{\alpha}_1,\symbf{\alpha}_2) & = -\max_{\symbf{\alpha}_1, \symbf{\alpha}_2} -g(\symbf{\alpha}_1,\symbf{\alpha}_2)                                                                                                                                             \\
	                                                                               & = -\max_{\symbf{\alpha}_1, \symbf{\alpha}_2} \min_{\symbf{w}, b, \symbf{\xi}} -\left(\frac{1}{2}\inner*{\symbf{\alpha}_1, \mathrm{Y} \mathrm{G} \mathrm{Y} \symbf{\alpha}_1} -  \inner*{\symbf{1}_M, \symbf{\alpha}_1} \right) \\
	                                                                               & = \min_{\symbf{\alpha}_1, \symbf{\alpha}_2} \left(\frac{1}{2}\inner*{\symbf{\alpha}_1, \mathrm{Y} \mathrm{G} \mathrm{Y} \symbf{\alpha}_1} -  \inner*{\symbf{1}_M, \symbf{\alpha}_1} \right)
\end{align*}

Then finally we can write the dual problem as:
\begin{align}\label{eq:dual-problem}
	\min_{\symbf{\alpha}} & \frac{1}{2} \inner*{\symbf{\alpha}, \mathrm{Y} \mathrm{G} \mathrm{Y} \symbf{\alpha}} -  \inner*{\symbf{1}_M, \symbf{\alpha}}
	\quad
	\text{subject to}
	\quad
	\begin{cases}
		0 \leq \symbf{\alpha} \leq C\symbf{1}_M \\
		\symbf{\alpha}^\top Y = 0
	\end{cases}\tag{DP}
\end{align}


% \section{Existence and Uniqueness of the Soft-Margin SVM Problem}
% \[
% \min_{w \in \mathbb{R}^d,\;b \in \mathbb{R},\;\xi \in \mathbb{R}^M}
% \quad \frac{1}{2}\,\|w\|^2 \;+\; C \sum_{i=1}^M \xi_i
% \quad\text{subject to}\quad
% \begin{cases}
% y_i\bigl(\langle w, x_i\rangle + b\bigr)\;\ge\;1 \;-\;\xi_i,\\[4pt]
% \xi_i \;\ge\; 0,
% \end{cases}
% \quad i = 1,\dots,M,
% \]

% \[
% \min_{\alpha \in \mathbb{R}^M} 
% \frac{1}{2} \langle \alpha, YGY \alpha \rangle 
% - \langle \mathbf{1}_M, \alpha \rangle
% \quad \text{subject to} \quad
% \begin{cases}
% \langle y, \alpha \rangle = 0,\\[6pt]
% 0 \le \alpha_i \le C
% \end{cases}
% \quad \text{for } i = 1, \dots, M.
% \]

% where each $y_i\in\{-1,+1\}$ and $C>0$ is a fixed penalty parameter.

% To prove the existence and uniqueness of a minimizer for a constrained optimization problem, 
% three properties are required. Feasibility ensures that the constraint set is non-empty, meaning 
% that there is at least one point that satisfies all constraints. Coercivity of the objective 
% function guarantees that the function tends to infinity as the norm of the variable tends to infinity,
% ensuring the existence of a global minimum on the feasible set. Finally, convexity of the objective 
% function implies that any local minimizer is also a global one, and if the function is strictly convex,
% it ensures that the minimizer is unique. Therefore, the combination of feasibility, coercivity, and 
% strict convexity, establishes both the existence and uniqueness of a global minimizer.


% \subsection*{1.\ Non-emptiness and Feasibility}

% The feasible set of $(w,b,\{\xi_i\})$ is non-empty. For example, one can always pick
% \[
% w = 0,\quad b = 0,\quad \xi_i = 1 \quad \forall\,i,
% \]
% so that
% \(
% y_i \langle w,x_i\rangle + b = 0
% \)
% and
% \(\,1 - \xi_i \le 0,\)
% thus satisfying
% $y_i(\langle w,x_i\rangle + b)\ge 1-\xi_i.$
% Hence there is at least one feasible point, guaranteeing the set of all feasible triples $(w,b,\xi)$ is non-empty.

% \subsection*{2.\ Coercivity and Existence of a Global Minimizer}

% Observe that the objective
% \[
% \frac{1}{2}\|w\|^2 \;+\; C\sum_{i=1}^M \xi_i
% \]
% is \emph{coercive} in the variables $(w,\xi)$: as $\|w\|\to\infty$ or $\sum_i \xi_i \to \infty,$ the cost grows unbounded. Together with the linear constraints $\xi_i\ge 0$ and the fact that the feasible set is not empty, standard results in convex analysis imply the existence of \emph{some} global minimizer $(w^*,b^*,\xi^*)$.

% \subsection*{3.\ Convexity and (Possible) Uniqueness}

% The function 
% \[
% \frac12 \|w\|^2 + C \sum_{i=1}^M \xi_i
% \]
% is convex in $(w,\xi)$, and \emph{affine} (hence convex) in $b$. Therefore, the entire objective is convex on a convex feasible region (defined by affine/linear inequalities). This guarantees that every local minimizer is also a global minimizer.

% \paragraph{Uniqueness.} 
% Strict convexity in the $w$-part implies that, \emph{for fixed $b$ and $\{\xi_i\}$}, the objective in $w$ is strictly convex, which typically ensures a unique $w^*$ \emph{conditional} on $b,\{\xi_i\}.$ However, because the objective is \emph{linear} in $b$, strict convexity does not extend to the $(b,w,\xi)$ space as a whole. In many cases, $b^*$ can still be pinned down by the margin constraints, but in some situations there may be multiple solutions that share the same minimal cost. 

% Hence, if the data are in ``general position,'' one typically obtains a \emph{unique} $w^*$ and a well-defined $b^*$, but strictly speaking, the problem is not globally \emph{strictly} convex in \emph{all} variables $(w,b,\xi)$ simultaneously. Therefore, uniqueness may or may not hold in full, depending on degeneracies in the data or constraints.


% \[
% f(\alpha) := \min_{\alpha \in \mathbb{R}^M} \frac{1}{2} \langle \alpha, YGY \alpha \rangle - \langle \mathbf{1}_M, \alpha \rangle \quad \text{s.t.} \quad 
% \begin{cases}
% \langle y, \alpha \rangle = 0 \\
% 0 \leq \alpha_i \leq C \quad \text{for } i = 1, \ldots, M
% \end{cases}
% \]

% \[
% Y = \text{diag}(y_1, \ldots, y_M) \quad \text{where } y_i = 1 \vee -1 \text{ for } i = 1, \ldots, m
% \]

% \[
% G = (\langle x_i, x_j \rangle)_{i,j=1,\ldots, M}, \quad \text{thus } G \text{ is the kernel matrix and SPD}
% \]

% \[
% \mathbf{1} = (1, \ldots, 1)^T
% \]

% \subsection*{1. Feasibility:}
% The area where the function $f(\alpha)$ can exist is given by
% \[
% \Omega = \left\{ \alpha \in \mathbb{R}^M \mid 0 \leq \alpha_i \leq C, \, \langle y, \alpha \rangle = 0 \right\}
% \]
% Thus the feasible region is the intersection of the area $[0, C]^M$ with the area $\sum_i y_i \alpha_i = 0$.

% Because there must at least be one positive and negative label, two indices, $p$ and $n$, can be picked such that $y_p = 1$ and $y_n = -1$ while $\alpha_p = \alpha_n$, letting all other $\alpha$ values be zero, satisfying both constraints. Guaranteeing that the feasible set is not zero.

% Also, $\alpha_n$ and $\alpha_p$ can be set to $0$ or $C$, thus boundary and corner points are also feasible.

% Thus $\Omega$ is a non-empty, closed and bounded area.

% \subsection*{2. Convexity:}
% Rewriting $f(\alpha) = \frac{1}{2} \langle \alpha, YGY \alpha \rangle - \langle \mathbf{1}_M, \alpha \rangle$ to

% \[
% f(\alpha) = \frac{1}{2} \alpha^T YGY \alpha - \mathbf{1}^T \alpha
% \]

% Finding the Hessian of this we find:

% \[
% \nabla f(\alpha) = Y G Y \alpha - \mathbf{1} \quad \Rightarrow \quad \nabla^2 f(\alpha) = Y G Y
% \]

% Where $G$ is semi-positive definite by definition as it is a gram matrix, and $Y$ is a diagonal matrix with $\pm 1$ values. Thus $YGY$ is symmetric and PSD, meaning $f$ is convex.



% \subsection{Lagrangian and KKT Conditions}

% The Lagrangian associated with the primal problem is given by:
% \[
% \mathcal{L}(\symbf{w}, b, \symbf{\xi}, \symbf{\alpha}, \symbf{\mu}) 
% = \frac{1}{2} \|\symbf{w}\|^2 + C \sum_{i=1}^M \xi_i
% - \sum_{i=1}^M \alpha_i \bigl( y_i (\langle \symbf{w}, \symbf{x}_i \rangle + b) - 1 + \xi_i \bigr)
% - \sum_{i=1}^M \mu_i \xi_i
% \]
% where $\symbf{\alpha} \ge 0$ and $\symbf{\mu} \ge 0$ are Lagrange multipliers for the margin and non-negativity constraints, respectively.

% From the KKT conditions, we obtain the following:

% \begin{itemize}
%   \item \textbf{Stationarity}:
%   \[
%   \nabla_{\symbf{w}} \mathcal{L} = \symbf{w} - \sum_{i=1}^M \alpha_i y_i \symbf{x}_i = 0 
%   \quad \Rightarrow \quad
%   \boxed{\symbf{w}^\star = \sum_{i=1}^M \alpha_i^\star y_i \symbf{x}_i}
%   \]
%   \[
%   \nabla_b \mathcal{L} = -\sum_{i=1}^M \alpha_i y_i = 0
%   \quad \Rightarrow \quad
%   \boxed{\sum_{i=1}^M \alpha_i^\star y_i = 0}
%   \]
%   \[
%   \nabla_{\xi_i} \mathcal{L} = C - \alpha_i - \mu_i = 0 
%   \quad \Rightarrow \quad
%   \boxed{\alpha_i^\star + \mu_i^\star = C}
%   \]

%   \item \textbf{Complementary slackness}:
%   \[
%   \alpha_i^\star \bigl( y_i(\langle \symbf{w}^\star, \symbf{x}_i \rangle + b^\star) - 1 + \xi_i^\star \bigr) = 0, 
%   \qquad
%   \mu_i^\star \xi_i^\star = 0
%   \]
% \end{itemize}

% \subsection{Recovering $b^\star$}

% For any support vector with $0 < \alpha_i^\star < C$, we know that $\xi_i^\star = 0$ and the corresponding constraint is active:
% \[
% y_i(\langle \symbf{w}^\star, \symbf{x}_i \rangle + b^\star) = 1 
% \quad \Rightarrow \quad 
% \boxed{b^\star = y_i - \langle \symbf{w}^\star, \symbf{x}_i \rangle}
% \]
% In practice, $b^\star$ is often computed by averaging this quantity over all such $i$ for which $0 < \alpha_i^\star < C$.

% \subsection{Recovering $\xi_i^\star$}

% Rewriting the margin constraint,
% \[
% \xi_i^\star \ge 1 - y_i(\langle \symbf{w}^\star, \symbf{x}_i \rangle + b^\star), \quad \xi_i^\star \ge 0,
% \]
% we get:
% \[
% \boxed{\xi_i^\star = \max\big(0,\; 1 - y_i(\langle \symbf{w}^\star, \symbf{x}_i \rangle + b^\star)\big)}
% \]

% \subsection{Summary}

% Given an optimal dual solution $\symbf{\alpha}^\star$, the corresponding primal solution is recovered via:
% \[
% \boxed{
% \begin{aligned}
% \symbf{w}^\star &= \sum_{i=1}^M \alpha_i^\star y_i \symbf{x}_i \\
% b^\star &= y_i - \langle \symbf{w}^\star, \symbf{x}_i \rangle \quad \text{for any } i \text{ such that } 0 < \alpha_i^\star < C \\
% \xi_i^\star &= \max\big(0,\; 1 - y_i(\langle \symbf{w}^\star, \symbf{x}_i \rangle + b^\star)\big)
% \end{aligned}
% }
% \]

% \begin{proof}{}{}
% 	\begin{align*}
% 		\min_{\symbf{w}, b, \symbf{\xi}} \quad
% 		 & f(\symbf{w}, b, \symbf{\xi})
% 		= \frac{1}{2}\|\symbf{w}\|_2^2 + C\,\symbf{1}_M^\top \symbf{\xi}
% 		\quad
% 		\text{subject to}
% 		\quad
% 		\begin{cases}
% 			\symbf{1}_M - Y\bigl(X\,\symbf{w} + b\,\symbf{1}_M\bigr) - \symbf{\xi} & \le \symbf{0}_M, \\[6pt]
% 			\symbf{\xi} \;\;                                                       & \ge \symbf{0}_M,
% 		\end{cases}\tag{P}
% 	\end{align*}
% 	\[
% 		\symbf{w}^\star = \symbf{y}^\top \symbf{\alpha}^\star \symbf{x} \tag{8}
% 	\]
% 	\[
% 		\symbf{b}^\star = \symbf{y} - {\symbf{w}^\star}^\top \symbf{x} = \symbf{y} - (\symbf{y}^\top \symbf{\alpha}^\star \symbf{x})^\top \symbf{x} = \symbf{y} - \symbf{x}^\top \symbf{\alpha}^{\star \top} \symbf{y} \symbf{x}
% 		= \symbf{y} - \inner*{\symbf{\alpha}^\star \symbf{x}, \symbf{y} \symbf{x}} \tag{9}
% 	\]

% 	\begin{align*}
% 		&\min_{\symbf{w}, b, \symbf{\xi}} \quad
% 		 f(\symbf{y}^\top \symbf{\alpha}^\star \symbf{x},\symbf{y} - {\symbf{w}^\star}^\top \symbf{x}, \symbf{\xi})
% 		= \frac{1}{2}\|\symbf{y}^\top \symbf{\alpha}^\star \symbf{x}\|_2^2 + C\,\symbf{1}_M^\top \symbf{\xi}
% 		\quad\\
% 		&\text{subject to}
% 		\quad
% 		\begin{cases}
% 			\symbf{1}_M - Y\bigl(X\,\symbf{y}^\top \symbf{\alpha}^\star \symbf{x} + (\symbf{y} - \inner*{\symbf{\alpha}^\star \symbf{x}, \symbf{y} \symbf{x}})\,\symbf{1}_M\bigr) - \symbf{\xi} & \le \symbf{0}_M, \\[6pt]
% 			\symbf{\xi} \;\;                                                                                     & \ge \symbf{0}_M,
% 		\end{cases}\\
% 	\end{align*}

% \end{proof}

% The dual of the primal problem~\eqref{eq:primal-problem} is obtained by forming the Lagrangian and eliminating $\symbf{w},b,\symbf{\xi}$:

% \begin{align}\label{eq:dual-problem}
% 	\min_{\symbf{\alpha}} 
% 	\;\;\frac{1}{2}\,\inner*{\symbf{\alpha},\,Y\,G\,Y\,\symbf{\alpha}} 
% 	\;-\;\inner*{\symbf{1}_M,\;\symbf{\alpha}}
% 	\quad
% 	\text{subject to}
% 	\quad
% 	\begin{cases}
% 		\inner{\symbf{y}}{\symbf{\alpha}} = 0,\\
% 		\symbf{0}_M \;\;\leq\;\; \symbf{\alpha} \;\;\leq\;\; C\,\symbf{1}_M.
% 	\end{cases}
% \end{align}

% Here $G = \bigl(\inner*{\symbf{x}_i}{\symbf{x}_j}\bigr)_{i,j=1}^M$ is the Gram matrix, and $Y=\operatorname{diag}(y_1,\dots,y_M)$.

% \begin{theorem}{Existence \& Uniqueness of Dual Solution}{existence-uniqueness-dual}
% The dual problem \eqref{eq:dual-problem} always admits at least one global minimizer. Moreover, if $YGY$ is strictly positive-definite on the subspace $\{\symbf{\alpha}\mid \inner{\symbf{y}}{\symbf{\alpha}}=0\}$, then its solution $\symbf{\alpha}^\star$ is unique. In degenerate cases, multiple $\symbf{\alpha}^\star$ can yield the same minimum.
% \end{theorem}

% \begin{proof}
% 	\paragraph{Existence of a solution.}
% 	The dual \eqref{eq:dual-problem} is a convex quadratic program in $\symbf{\alpha}$, subject to box constraints $0 \le \alpha_i \le C$ and the linear equality $\inner{\symbf{y}}{\symbf{\alpha}}=0$. Thus:
% 	\begin{itemize}
% 		\item The feasible region (a slice of a hyperbox) is closed and bounded.
% 		\item The dual objective 
% 		\(
% 		  \tfrac12\, \symbf{\alpha}^\top (YGY)\,\symbf{\alpha}\;-\;\symbf{1}_M^\top \symbf{\alpha}
% 		\)
% 		is continuous and convex over that compact set.
% 	\end{itemize}
% 	By the extreme-value theorem, a global minimizer $\symbf{\alpha}^\star$ must exist.
% 	\paragraph{Uniqueness of a solution.}
% 	\begin{itemize}
% 		\item If $Y\,G\,Y$ is positive-definite on the subspace $\{\symbf{\alpha}\mid \inner{\symbf{y}}{\symbf{\alpha}}=0\}$, then the quadratic objective is strictly convex there, guaranteeing a single minimizer.
% 		\item In degenerate situations, e.g.\ if $G$ is positive semidefinite but not definite on that subspace, there can be infinitely many optimal $\symbf{\alpha}^\star$.
% 	\end{itemize}
% \end{proof}



\footnote{The Lagrangian is a function that combines the objective and constraints of a 
constrained optimization problem. For a problem of the form
\[
\min_{x \in \mathbb{R}^n} \; f(x) \quad \text{subject to} \quad g_i(x) \le 0 \; \text{for } i = 1, 
\ldots, m, \quad h_j(x) = 0 \; \text{for } j = 1, \ldots, p,
\]
the Lagrangian is defined as
\[
\mathcal{L}(x, \boldsymbol{\lambda}, \boldsymbol{\mu}) 
= f(x) + \sum_{i=1}^m \lambda_i g_i(x) + \sum_{j=1}^p \mu_j h_j(x) 
= f(x) + \boldsymbol{\lambda}^\top g(x) + \boldsymbol{\mu}^\top h(x),
\]
And $g(x) = [g_1(x), \ldots, g_m(x)]^\top$ and $h(x) = [h_1(x), \ldots, h_p(x)]^\top$.
}

\footnote{The Karush-Kuhn-Tucker (KKT) conditions are necessary conditions 
for a point $x^\star$ to be optimal in constrained optimization, and sufficient when the problem 
is convex. These conditions are Stationarity, Primal and Dual feasibility, and Complementary slackness. 
Stationarity constitutes that the gradient of the Lagrangian with respect to $x$ vanishes at $x^\star$. 
Primal feasibility is when the original constraints are satisfied. Dual feasibility requires 
$\lambda_i \ge 0$ for all $i$. Complementary slackness is that each inequality constraint satisfies 
$\lambda_i g_i(x^\star) = 0$.
}